\documentclass[11pt,a4paper]{article}

\usepackage[margin=1in]{geometry}
\usepackage{booktabs}
\usepackage{longtable}
\usepackage{array}
\usepackage{multirow}
\usepackage{xcolor}
\usepackage{colortbl}
\usepackage{hyperref}
\usepackage{enumitem}
\usepackage{titlesec}
\usepackage{fancyhdr}
\usepackage{graphicx}
\usepackage{tabularx}
\usepackage{makecell}
\usepackage{mdframed}
\usepackage{amsmath}
\usepackage{fontenc}
\usepackage[utf8]{inputenc}
\usepackage{lmodern}
\usepackage{microtype}

% Colors
\definecolor{phaseblue}{HTML}{1F4E79}
\definecolor{lightgray}{HTML}{F2F2F2}
\definecolor{darkgray}{HTML}{404040}
\definecolor{accentblue}{HTML}{2E75B6}

% Hyperref setup
\hypersetup{
    colorlinks=true,
    linkcolor=accentblue,
    urlcolor=accentblue,
    citecolor=accentblue
}

% Header/Footer
\pagestyle{fancy}
\fancyhf{}
\rhead{\textcolor{darkgray}{\small Multilingual Image Editing Error Classifier}}
\lhead{\textcolor{darkgray}{\small Open-Source Model Benchmark}}
\cfoot{\thepage}
\renewcommand{\headrulewidth}{0.4pt}

% Section formatting
\titleformat{\section}{\large\bfseries\color{phaseblue}}{}{0em}{}[\titlerule]
\titleformat{\subsection}{\normalsize\bfseries\color{accentblue}}{}{0em}{}

% Phase box style
\newmdenv[
  linecolor=phaseblue,
  linewidth=2pt,
  backgroundcolor=lightgray,
  innertopmargin=8pt,
  innerbottommargin=8pt,
  innerleftmargin=10pt,
  innerrightmargin=10pt,
  skipabove=10pt,
  skipbelow=10pt
]{phasebox}

% Highlight box
\newmdenv[
  linecolor=accentblue,
  linewidth=1pt,
  backgroundcolor=lightgray,
  innertopmargin=6pt,
  innerbottommargin=6pt,
  innerleftmargin=8pt,
  innerrightmargin=8pt,
  skipabove=6pt,
  skipbelow=6pt
]{infobox}

\begin{document}

% Title
\begin{center}
    {\LARGE\bfseries Multilingual Image Editing Error Classifier \&}\\[4pt]
    {\LARGE\bfseries Open-Source Model Benchmark}\\[10pt]
    {\normalsize\itshape Based on MagicBrush Dataset \quad|\quad Nepali $\cdot$ Bangla $\cdot$ Hindi Translation + Error Annotation}\\[14pt]

    % Phase pipeline bar
    \begin{tabular}{|>{\centering\arraybackslash\columncolor{phaseblue}\color{white}}p{3cm}|>{\centering\arraybackslash\columncolor{accentblue!60}}p{3cm}|>{\centering\arraybackslash\columncolor{accentblue!60}}p{3cm}|>{\centering\arraybackslash\columncolor{accentblue!60}}p{3cm}|}
        \hline
        \textbf{Phase 1} \\ Data Prep \& Translation &
        \textbf{Phase 2} \\ Annotation &
        \textbf{Phase 3} \\ Model Training &
        \textbf{Phase 4} \\ Benchmarking \\
        \hline
    \end{tabular}
\end{center}

\vspace{10pt}

%% ============================================================
\section{1. Project Overview}
%% ============================================================

This plan outlines a four-phase research pipeline that begins with the MagicBrush dataset, creates a multilingual annotated benchmark, trains a custom error classifier, and uses that classifier to benchmark open-source image editing models.

\subsection*{Dataset Source}

\begin{infobox}
\texttt{osunlp/MagicBrush}\\
\url{https://huggingface.co/datasets/osunlp/MagicBrush}
\end{infobox}

\vspace{6pt}
\begin{tabularx}{\textwidth}{>{\bfseries}lX}
\toprule
source\_img & Input image — original real image (turn 1) or edited image from previous turn (turn $\geq$ 2) \\
\midrule
instruction & Natural language instruction describing the desired edit (English) \\
\midrule
target\_img & Output image after the editing model processes the instruction \\
\bottomrule
\end{tabularx}

\subsection{Research Goals}

\begin{tabularx}{\textwidth}{>{\centering\bfseries\arraybackslash\columncolor{phaseblue}\color{white}}p{1.5cm}X}
\toprule
Goal 1 & Translate \texttt{instruction} into Nepali, Bangla, and Hindi using a multilingual LLM pipeline \\[4pt]
\rowcolor{lightgray}
\textbf{Goal 2} & Annotate each (\texttt{source\_img}, \texttt{instruction}, \texttt{target\_img}) triple with error labels from a defined 11-class taxonomy \\[4pt]
Goal 3 & Train a multimodal error classifier on the annotated dataset \\[4pt]
\rowcolor{lightgray}
\textbf{Goal 4} & Evaluate multiple open-source image editing models using the trained classifier as the judge \\[4pt]
Goal 5 & Publish a benchmark leaderboard comparing model performance across error categories and languages \\
\bottomrule
\end{tabularx}

%% ============================================================
\section{2. Phase 1 — Data Preparation \& Translation}
%% ============================================================

\begin{phasebox}
\textbf{PHASE 1} \quad Data Preparation \& Multilingual Translation\\
\textit{Subset selection, cleaning, and LLM-based prompt translation}
\end{phasebox}

\subsection{2.1 Dataset Sampling Strategy}

The MagicBrush dataset contains 10,308 edit turns across train and dev splits (8,807 train + 528 dev + held-out test). It covers single-turn and multi-turn editing scenarios. For initial experiments, a stratified sample is used to ensure edit diversity. The recommended subset sizes are:

\vspace{6pt}
\begin{tabularx}{\textwidth}{>{\bfseries}lX}
\toprule
Annotation Pilot & 500 samples — diverse edit types, manual quality check \\
\midrule
Full Annotation Set & 5,000–10,000 samples — balanced across edit categories \\
\midrule
Translation Set & All 5,000–10,000 samples translated into 3 languages \\
\midrule
Model Training Set & $\sim$8,000 annotated + translated samples (80\% train / 20\% val) \\
\bottomrule
\end{tabularx}

\subsection{2.2 Translation Pipeline}

Each \texttt{instruction} will be translated from English into three low-resource South Asian languages using the following approach:

\begin{infobox}
\textbf{Translation Approach: IndicTrans2 + GPT-4 Fallback}\\[4pt]
\textbf{Primary:} IndicTrans2 (AI4Bharat) — open-source, supports Nepali/Bangla/Hindi natively\\
\textbf{Fallback:} GPT-4o or Google Cloud Translate API for quality verification\\
\textbf{Validation:} Back-translation (Nepali$\to$English, compare with original using BLEU/chrF)
\end{infobox}

\subsubsection*{Translation Quality Control}

\begin{itemize}[leftmargin=1.5em]
    \item All translated prompts back-translated to English and similarity scored (BLEU $\geq$ 0.65 threshold)
    \item Low-confidence translations flagged for manual review by a native speaker
    \item Culturally ambiguous terms preserved in original English with annotation note
    \item Final schema adds three new columns: \texttt{instruction\_ne}, \texttt{instruction\_bn}, \texttt{instruction\_hi}
\end{itemize}

\subsection{2.3 Output Schema After Phase 1}

\begin{tabularx}{\textwidth}{>{\ttfamily\bfseries}lX}
\toprule
source\_img & PIL Image / path — input image (original or from previous edit turn) \\
\rowcolor{lightgray}
instruction & str — original English editing instruction \\
instruction\_ne & str — Nepali translation of instruction \\
\rowcolor{lightgray}
instruction\_bn & str — Bangla translation of instruction \\
instruction\_hi & str — Hindi translation of instruction \\
\rowcolor{lightgray}
target\_img & PIL Image / path — model-generated edited image \\
translation\_confidence & float[3] — back-translation BLEU scores per language \\
\bottomrule
\end{tabularx}

%% ============================================================
\section{3. Phase 2 — Error Annotation}
%% ============================================================

\begin{phasebox}
\textbf{PHASE 2} \quad Error Annotation\\
\textit{Applying the 11-class error taxonomy to each image editing triple}
\end{phasebox}

\subsection{3.1 Error Taxonomy (11 Classes)}

Each triple (\texttt{source\_img}, \texttt{instruction}, \texttt{target\_img}) will be labeled with one or more error types from the following taxonomy. Labels are \textit{not mutually exclusive} — a single edited image can have multiple errors.

\vspace{6pt}
\begin{tabularx}{\textwidth}{>{\centering\bfseries}p{0.5cm} >{\bfseries}p{3.2cm} X >{\itshape}p{4cm}}
\toprule
\textbf{\#} & \textbf{Error Type} & \textbf{Description} & \textbf{Example} \\
\midrule
0 & Wrong Object & The edited image contains the wrong object type & Asked for cat, got dog \\
\rowcolor{lightgray}
1 & Missing Object & A required object is absent from the output & Remove table — chair also gone \\
2 & Extra Object & Unwanted additional objects appear & Add lamp — random couch added \\
\rowcolor{lightgray}
3 & Wrong Attribute & Object exists but with incorrect color/size/shape & Make sky red — sky turns green \\
4 & Spatial Error & Objects are mispositioned or wrongly oriented & Move left — object moves right \\
\rowcolor{lightgray}
5 & Style Mismatch & Editing style conflicts with original image style & Anime filter on photorealistic image \\
6 & Over-editing & Too many regions changed beyond the instruction & Change hat — face also changed \\
\rowcolor{lightgray}
7 & Under-editing & The instruction is only partially applied & Make background dark — only half dark \\
8 & Artifact / Quality & Noise, blur, distortion, or rendering errors & Visible seams, blurry patches \\
\rowcolor{lightgray}
9 & Ambiguous Prompt & Prompt is unclear so model guesses incorrectly & ``Make it better'' — undefined \\
10 & Failed Removal & Object still present after removal instruction & Remove tree — tree still visible \\
\bottomrule
\end{tabularx}

\subsection{3.2 Annotation Strategy}

\subsubsection*{Human Annotation}
\begin{itemize}[leftmargin=1.5em]
    \item \textbf{Platform:} Label Studio (self-hosted) or Scale AI
    \item \textbf{Annotators see:} original image + prompt + edited image side by side
    \item \textbf{Task:} select all applicable error labels (multi-label), 0 = no error
    \item \textbf{Inter-Annotator Agreement:} minimum 2 annotators per sample, Cohen's Kappa $\geq$ 0.7
    \item \textbf{Estimated effort:} $\sim$1.5 min/sample $\times$ 5,000 = $\sim$125 hours total annotation time
\end{itemize}

\subsubsection*{Semi-Automated Annotation}
\begin{itemize}[leftmargin=1.5em]
    \item Use GPT-4V / LLaVA with structured prompts to pre-label all samples
    \item Human annotators only review low-confidence predictions ($<$ 0.7 model confidence)
    \item Reduces human effort by $\sim$60\% while maintaining quality on 5,000+ samples
\end{itemize}

\begin{infobox}
\textbf{Hybrid Approach}\\[4pt]
\textbf{Step 1:} GPT-4V auto-labels all 5,000 samples with confidence scores\\
\textbf{Step 2:} Human annotators review top 20\% lowest-confidence samples ($\sim$1,000)\\
\textbf{Step 3:} Agreement check on 200-sample validation set (2 humans + model)\\
\textbf{Step 4:} Final annotation stored as multi-hot vector: e.g.\ \texttt{[0,0,0,1,0,0,1,0,0,0,0]}
\end{infobox}

\subsection{3.3 Final Annotation Schema}

\begin{tabularx}{\textwidth}{>{\ttfamily\bfseries}lX}
\toprule
error\_label\_vector & int[11] — multi-hot, e.g.\ [0,0,1,1,0,0,0,0,0,0,0] \\
\rowcolor{lightgray}
error\_types & List[str] — human-readable error names, e.g.\ ['Extra Object', 'Wrong Attribute'] \\
annotation\_confidence & float — model confidence or human agreement score \\
\rowcolor{lightgray}
annotator\_id & str — annotator identifier for traceability \\
has\_error & bool — True if any error exists (for binary classification baseline) \\
\bottomrule
\end{tabularx}

%% ============================================================
\section{4. Phase 3 — Model Training (Error Classifier)}
%% ============================================================

\begin{phasebox}
\textbf{PHASE 3} \quad Error Classifier Training\\
\textit{Fine-tuning a multimodal model on annotated triplets}
\end{phasebox}

\subsection{4.1 Task Formulation}

The error classifier takes as input the full editing triple and outputs a multi-label prediction:

\begin{infobox}
\textbf{Input:} \quad \texttt{source\_img} + \texttt{instruction} (any language) + \texttt{target\_img}\\[4pt]
\textbf{Output:} \quad 11-dim multi-hot vector $\to$ \texttt{[0, 1, 0, 0, 0, 0, 1, 0, 0, 0, 0]}
\end{infobox}

\begin{infobox}
\textbf{Architecture}\\[4pt]
BLIP-2 (frozen encoder) + XLM-RoBERTa (for multilingual text) + dual-image attention + 2-layer MLP classifier\\[4pt]
\textit{Why: Supports multilingual prompts naturally, image-pair comparison baked in, runs on single A100 (40GB)}
\end{infobox}

\subsection{4.3 Training Configuration}

\begin{tabularx}{\textwidth}{>{\bfseries}lX}
\toprule
Loss Function & Binary Cross-Entropy with Logits (multi-label, class-weighted) \\
\rowcolor{lightgray}
Optimizer & AdamW with cosine LR schedule, warmup 500 steps \\
Learning Rate & 1e-4 (MLP head), 1e-5 (encoder, if unfrozen) \\
\rowcolor{lightgray}
Batch Size & 16 (gradient accumulation $\times$ 4 = effective 64) \\
Epochs & 10–20 with early stopping (patience = 3 on val F1) \\
\rowcolor{lightgray}
Data Split & 80\% train / 10\% validation / 10\% test (stratified) \\
Class Weights & Inverse frequency weighting to handle class imbalance \\
\rowcolor{lightgray}
Augmentation & Random crop, color jitter on \texttt{source\_img} only (not \texttt{target\_img}) \\
\bottomrule
\end{tabularx}

\subsection{4.4 Evaluation Metrics}

\begin{tabularx}{\textwidth}{>{\bfseries}lX}
\toprule
Per-class F1 & Precision, Recall, F1 for each of the 11 error types \\
\rowcolor{lightgray}
Macro F1 & Average F1 across all classes (handles imbalance) \\
Weighted F1 & F1 weighted by class frequency — \textbf{primary metric} \\
\rowcolor{lightgray}
Hamming Loss & Fraction of incorrectly predicted labels \\
Exact Match & \% samples where all 11 labels predicted correctly \\
\rowcolor{lightgray}
AUC-ROC & Per-class area under ROC curve \\
\bottomrule
\end{tabularx}

%% ============================================================
\section{5. Phase 4 — Benchmarking Open-Source Models}
%% ============================================================

\begin{phasebox}
\textbf{PHASE 4} \quad Open-Source Model Benchmarking\\
\textit{Generate images with open-source models and evaluate with our classifier}
\end{phasebox}

\subsection{5.1 Open-Source Models to Benchmark}

\begin{tabularx}{\textwidth}{>{\bfseries}lX}
\toprule
InstructPix2Pix (SD 1.5) & \texttt{timbrooks/instruct-pix2pix} — Baseline model, most widely cited \\
\rowcolor{lightgray}
MagicBrush Fine-tuned & \texttt{osunlp/InstructPix2Pix-MagicBrush} — Human-annotated fine-tune \\
UltraEdit (SD3) & \texttt{pkunlp-icler/UltraEdit} — NeurIPS 2024 SOTA, 4M training samples \\
\rowcolor{lightgray}
MGIE (LLaVA-guided) & \texttt{tsujuifu/mllm-mgie} — Uses LLM to rewrite prompts before editing \\
HIVE & \texttt{zhanglingao/hive} — Human-feedback aligned editing model \\
\rowcolor{lightgray}
Your Fine-tuned Model & Custom model trained on translated + annotated dataset \\
\bottomrule
\end{tabularx}

\subsection{5.2 Evaluation Pipeline}

\begin{infobox}
\textbf{Benchmarking Pipeline Steps}\\[4pt]
\textbf{Input:} hold-out test set of 500–1,000 \texttt{source\_img} + \texttt{instruction} triples (unseen during training)\\[4pt]
For each open-source model, generate edited images using English, Nepali, Bangla, and Hindi prompts separately\\[4pt]
Run each (\texttt{source\_img}, \texttt{instruction}, \texttt{generated\_target\_img}) triple through your trained classifier\\[4pt]
Record error type predictions and confidence scores per model per language\\[4pt]
Compute automatic metrics: CLIP-I similarity, DINO score, FID (where ground truth available)\\[4pt]
Aggregate results into a benchmark leaderboard table
\end{infobox}

\subsection{5.3 Benchmark Metrics}

\begin{tabularx}{\textwidth}{>{\bfseries}lX}
\toprule
Error Rate ($\downarrow$) & \% of generated images flagged with at least 1 error by classifier \\
\rowcolor{lightgray}
Per-Error Distribution & Which error types each model is prone to (error profile) \\
CLIP Image Similarity ($\uparrow$) & Cosine similarity between original and edited image CLIP embeddings \\
\rowcolor{lightgray}
DINO Score ($\uparrow$) & Structure-preserving metric; how well model follows edit without destroying image \\
FID ($\downarrow$) & Fr\'{e}chet Inception Distance — overall image quality vs ground truth (if available) \\
\rowcolor{lightgray}
Cross-Lingual Gap ($\downarrow$) & Difference in error rate between English and Nepali/Bangla/Hindi prompts \\
\bottomrule
\end{tabularx}

\subsection{5.5 Cross-Lingual Analysis}

A key contribution of this work is evaluating how well open-source models handle low-resource language prompts. The cross-lingual gap metric measures the degradation in editing quality when using Nepali/Bangla/Hindi versus English:

\vspace{6pt}
\begin{tabularx}{\textwidth}{>{\bfseries}lX}
\toprule
English (Baseline) & \\
\rowcolor{lightgray}
Nepali & \\
Bangla & \\
\rowcolor{lightgray}
Hindi & \\
\bottomrule
\end{tabularx}

%% ============================================================
\section{6. Tools, Libraries \& Infrastructure}
%% ============================================================

\subsection{6.1 Core Libraries}

\begin{tabularx}{\textwidth}{>{\bfseries}lX}
\toprule
Dataset Loading & \texttt{datasets} (HuggingFace), \texttt{Pillow}, \texttt{pandas} \\
\rowcolor{lightgray}
Translation & \texttt{IndicTrans2} (AI4Bharat), \texttt{sacrebleu} (BLEU scoring) \\
Annotation UI & Label Studio (self-hosted), \texttt{argilla} \\
\rowcolor{lightgray}
Model Training & \texttt{transformers}, \texttt{PEFT/LoRA}, \texttt{accelerate}, \texttt{bitsandbytes} \\
Image Metrics & \texttt{torch-fidelity} (FID), \texttt{open\_clip} (CLIP-I), \texttt{timm} (DINO) \\
\rowcolor{lightgray}
Experiment Tracking & Weights \& Biases (\texttt{wandb}) or MLflow \\
Inference & \texttt{diffusers} (for open-source editing models) \\
\bottomrule
\end{tabularx}

\end{document}
